\documentclass[11pt,a4paper]{article}

\usepackage[T1]{fontenc}
\usepackage[utf8]{inputenc}
\usepackage{amsmath,amssymb,amsthm}
\usepackage[margin=2.6cm]{geometry}
\usepackage{hyperref}

\newtheorem{theorem}{Theorem}[section]
\newtheorem{proposition}[theorem]{Proposition}
\newtheorem{lemma}[theorem]{Lemma}
\newtheorem{corollary}[theorem]{Corollary}
\theoremstyle{definition}
\newtheorem{definition}[theorem]{Definition}
\newtheorem{remark}[theorem]{Remark}

\newcommand{\F}{\mathbb{F}}
\newcommand{\GF}{\mathrm{GF}}
\newcommand{\Tr}{\operatorname{Tr}}
\newcommand{\rhs}{R}

\title{The Muller--Cohen--Matthews polynomial is a permutation polynomial\\
\large A human-readable account of the formalised proof}
\author{}
\date{}

\begin{document}
\maketitle

\begin{abstract}
We give a complete, human-readable proof of the following statement: for positive
integers $k<n$ with $\gcd(k,n)=1$ and $k$ odd, the Muller--Cohen--Matthews (MCM)
polynomial
\[
  M_{n,k}(z)\;=\;z\Bigl(\sum_{i=0}^{k-1}z^{2^{i}-1}\Bigr)^{2^{k}+1}
\]
is a permutation polynomial of $L=\GF(2^{n})$.  The proof follows Dobbertin's
route: one first proves the permutation criterion for the generalized Kasami
polynomials $q_\alpha$, and then transports it to the MCM side through the pair of
mutually inverse linearized transfer maps $\psi_\beta$, $\varphi_\alpha$
(Dobbertin's Theorem 4).  Every step below corresponds to a machine-checked Lean
statement of the accompanying development; the correspondence is tabulated in
Appendix~\ref{app:lean}.
\end{abstract}

\tableofcontents

\section{Setting and notation}

Throughout, $n$ and $k$ are positive integers with
\[
  k<n,\qquad \gcd(k,n)=1 ,
\]
and $L=\GF(2^{n})$ is the field with $2^{n}$ elements, of characteristic two.  We
write $k'$ for a positive representative of the inverse of $k$ modulo $n$, i.e.
\[
  k\,k'\equiv 1 \pmod n ,\qquad 0<k'<n .
\]
Since $\gcd(k,n)=1$ such a $k'$ exists, and together with it we fix the Bézout
quotient $n'\in\mathbb{N}$ determined by
\begin{equation}\label{eq:bezout}
  k'k \;=\; n'n+1 .
\end{equation}

Two standing facts about $L$ are used constantly and without further comment.

\begin{lemma}\label{lem:frob}
For all $x\in L$ we have $x^{2^{n}}=x$; consequently $x^{2^{r}}$ depends only on
$r\bmod n$.  Moreover $u\mapsto u^{2^{r}}$ is a field automorphism, so
$(u+v)^{2^{r}}=u^{2^{r}}+v^{2^{r}}$, and $-1=1$.
\end{lemma}

\begin{lemma}\label{lem:gcdpow}
For all $r\ge 1$, $\gcd(2^{r}-1,2^{n}-1)=2^{\gcd(r,n)}-1$ and
$\gcd(2^{r},2^{n}-1)=1$.  In particular, since $\gcd(k,n)=1$:
\begin{enumerate}
\item $u\mapsto u^{2^{k}-1}$ is a bijection of $L^{\times}$;
\item $u\mapsto u^{2^{k}}$ is a bijection of $L$, and $u^{2^{k}}=1$ forces $u=1$;
\item the kernel of the additive map $u\mapsto u^{2^{k}}+u$ is
      $\F_{2^{\gcd(k,n)}}=\{0,1\}$.
\end{enumerate}
\end{lemma}

The \emph{absolute trace polynomial} is
\[
  \Tr(z)\;=\;\sum_{i=0}^{n-1}z^{2^{i}} ,
\]
an additive map $L\to\F_2\subseteq L$ satisfying $\Tr(z)^{2^{r}}=\Tr(z)$ for all
$r$.

\begin{definition}[Generalized Kasami polynomial]
For $\alpha\in\{0,1\}$ put
\[
  q_\alpha(z)\;=\;\Bigl(\sum_{i=1}^{k'}z^{2^{ik}}+\alpha\Tr(z)\Bigr)\,
                 z^{(2^{n}-1)-(2^{k}+1)} .
\]
\end{definition}

\begin{definition}[Generalized MCM polynomial]
For $\beta\in\{0,1\}$ put
\[
  P_\beta(z)\;=\;\Bigl(\sum_{i=0}^{k-1}z^{2^{i}}+\beta\Tr(z)\Bigr)^{2^{k}+1}
                 z^{(2^{n}-1)-2^{k}} .
\]
\end{definition}

\begin{definition}[MCM polynomial]
\[
  M_{n,k}(z)\;=\;z\Bigl(\sum_{i=0}^{k-1}z^{2^{i}-1}\Bigr)^{2^{k}+1}.
\]
\end{definition}

In each case we distinguish carefully between the \emph{polynomial} and the
\emph{evaluation map} $L\to L$ it induces; ``$p$ is a permutation polynomial''
means that the evaluation map $x\mapsto p(x)$ is a bijection of $L$.  As usual we
use the convention $0^{-1}=0$, so that $x\mapsto x^{-1}$ is a bijection of $L$
fixing $0$.

The goal of this note is:

\begin{theorem}[Main theorem]\label{thm:main}
Let $0<k<n$ with $\gcd(k,n)=1$ and $k$ odd.  Then $M_{n,k}$ is a permutation
polynomial of $\GF(2^{n})$.
\end{theorem}

The proof occupies Sections \ref{sec:kasami}--\ref{sec:mcm}: Section
\ref{sec:kasami} proves the permutation criterion for $q_\alpha$, Section
\ref{sec:transfer} the transfer theorem relating $q_\alpha$ and $P_\beta$, and
Section \ref{sec:mcm} deduces Theorem \ref{thm:main}.

\section{The permutation criterion for the generalized Kasami polynomials}
\label{sec:kasami}

The aim of this section is the following criterion, which is where all the work
lies.

\begin{theorem}\label{thm:kasami}
Let $0<k<n$, $\gcd(k,n)=1$, $0<k'<n$ with $kk'\equiv1\pmod n$, and
$\alpha\in\{0,1\}$.  Then $q_\alpha$ is a permutation polynomial of $\GF(2^{n})$
if and only if
\[
  k'+\alpha n\equiv 1 \pmod 2 .
\]
\end{theorem}

Throughout this section we abbreviate
\[
  S(z)\;=\;\sum_{i=1}^{k'}z^{2^{ik}},\qquad
  \rhs(z)\;=\;S(z)+\alpha\Tr(z),
\]
so that $q_\alpha(z)=\rhs(z)\,z^{(2^{n}-1)-(2^{k}+1)}$.  Both $S$ and $\rhs$ are
additive maps $L\to L$.

\subsection{Fibres of $q_\alpha$ as a Kasami equation}

\begin{lemma}\label{lem:fibre}
Let $c,z\in L$.
\begin{enumerate}
\item If $z\neq0$ then $q_\alpha(z)=c$ if and only if
      \begin{equation}\label{eq:K}
        c\,z^{2^{k}+1}\;=\;\rhs(z). \tag{K}
      \end{equation}
\item $q_\alpha(z)=0$ if and only if $\rhs(z)=0$ (including for $z=0$).
\end{enumerate}
\end{lemma}

\begin{proof}
Assume $z\neq0$, so $z^{2^{n}-1}=1$.  Since $k<n$ we have $2^{k}+1\le 2^{n}-1$,
hence multiplying $q_\alpha(z)=\rhs(z)z^{(2^{n}-1)-(2^{k}+1)}$ by $z^{2^{k}+1}$
gives $q_\alpha(z)z^{2^{k}+1}=\rhs(z)z^{2^{n}-1}=\rhs(z)$.  As $z^{2^{k}+1}\neq0$,
$q_\alpha(z)=c$ is equivalent to \eqref{eq:K}.  This proves (1), and (2) for
$z\neq 0$ (take $c=0$).  For $z=0$ both $q_\alpha(0)$ and $\rhs(0)$ vanish, because
every exponent occurring in $S$ and $\Tr$ is positive.
\end{proof}

\subsection{Necessity of the parity condition}

\begin{proposition}\label{prop:necessity}
If $q_\alpha$ is a permutation polynomial then $k'+\alpha n\equiv1\pmod2$.
\end{proposition}

\begin{proof}
Suppose $k'+\alpha n$ is even.  Then
\[
  \rhs(1)=\sum_{i=1}^{k'}1+\alpha\sum_{i=0}^{n-1}1=(k'+\alpha n)\cdot 1=0
\]
in $L$, because $L$ has characteristic two.  By Lemma \ref{lem:fibre}(2),
$q_\alpha(1)=0=q_\alpha(0)$, while $1\neq0$ in $L$.  Hence $q_\alpha$ is not
injective.
\end{proof}

\subsection{Frobenius identities and the linearized equation}

\begin{lemma}\label{lem:shift}
For every $x\in L$,
\[
  S(x)^{2^{k}}\;=\;S(x)+x^{2^{k}}+x^{2^{k+1}},\qquad
  \Tr(x)^{2^{k}}=\Tr(x),
\]
and consequently $\rhs(x)^{2^{k}}=\rhs(x)+x^{2^{k}}+x^{2^{k+1}}$.
\end{lemma}

\begin{proof}
Applying the automorphism $u\mapsto u^{2^{k}}$ termwise,
$S(x)^{2^{k}}=\sum_{i=1}^{k'}x^{2^{(i+1)k}}=\sum_{i=2}^{k'+1}x^{2^{ik}}$.  Hence,
in characteristic two,
\[
  S(x)^{2^{k}}=S(x)+x^{2^{k}}+x^{2^{(k'+1)k}} .
\]
By \eqref{eq:bezout}, $(k'+1)k=k'k+k\equiv 1+k \pmod n$, so by Lemma
\ref{lem:frob} $x^{2^{(k'+1)k}}=x^{2^{k+1}}$.  The trace identity is the standard
invariance of $\Tr$ under Frobenius, and the last claim follows since
$\alpha\in\{0,1\}$ satisfies $\alpha^{2^{k}}=\alpha$.
\end{proof}

\begin{definition}
For $c\in L$ let
\[
  \ell_c(X)=c^{2^{k}}X^{2^{2k}}+X^{2^{k}}+cX+1,\qquad
  \ell_{0,c}(X)=\ell_c(X)+1=c^{2^{k}}X^{2^{2k}}+X^{2^{k}}+cX .
\]
$\ell_{0,c}$ is $\F_2$-linear, so $\ell_c(x)=\ell_c(y)$ if and only if
$\ell_{0,c}(x+y)=0$.
\end{definition}

\begin{proposition}\label{prop:linearized}
If $x\neq0$ satisfies \eqref{eq:K}, then $\ell_c(x)=0$.
\end{proposition}

\begin{proof}
Raise \eqref{eq:K} to the power $2^{k}$ and add \eqref{eq:K} to it.  Using Lemma
\ref{lem:shift} on the right-hand side,
\[
  c^{2^{k}}x^{2^{2k}+2^{k}}+c\,x^{2^{k}+1}
  \;=\;\rhs(x)^{2^{k}}+\rhs(x)\;=\;x^{2^{k}}+x^{2^{k+1}} .
\]
Dividing by $x^{2^{k}}\neq0$ and using $x^{2^{k+1}}/x^{2^{k}}=x^{2^{k}}$ yields
\[
  c^{2^{k}}x^{2^{2k}}+c\,x+1+x^{2^{k}}=0,
\]
which is $\ell_c(x)=0$.
\end{proof}

\begin{proposition}\label{prop:zerofibre}
Assume $k'+\alpha n\equiv 1 \pmod 2$.  If $\rhs(x)=0$ then $x=0$.
\end{proposition}

\begin{proof}
From $\rhs(x)=0$ and Lemma \ref{lem:shift},
$0=\rhs(x)^{2^{k}}+\rhs(x)=x^{2^{k}}+x^{2^{k+1}}=x^{2^{k}}\bigl(1+x^{2^{k}}\bigr)$.
Hence $x=0$ or $x^{2^{k}}=1$; in the latter case $x=1$ by Lemma
\ref{lem:gcdpow}(2).  But $\rhs(1)=(k'+\alpha n)\cdot1=1\neq0$ by the parity
hypothesis.  So $x=0$.
\end{proof}

\subsection{First case: $c$ is not of the form $\gamma^{2^{k}+1}+\gamma$}

\begin{lemma}[Factorisation]\label{lem:factorization}
Let $c\neq0$ and $z\in L$.  Put $D=c\,z^{2^{k}-1}$ and $A=D^{2^{n-1}}$, so that
$A^{2}=D$.  Then
\[
  \ell_{0,c}(z)\;=\;c^{-1}\bigl(A^{2^{k}+1}+A+c\bigr)^{2}\,z .
\]
\end{lemma}

\begin{proof}
Squaring is additive, so
$\bigl(A^{2^{k}+1}+A+c\bigr)^{2}=A^{2(2^{k}+1)}+A^{2}+c^{2}
 =D^{2^{k}+1}+D+c^{2}$.  Now
\[
  c^{-1}zD^{2^{k}+1}=c^{-1}z\,c^{2^{k}+1}z^{(2^{k}-1)(2^{k}+1)}
   =c^{2^{k}}z^{2^{2k}},\qquad
  c^{-1}zD=c^{-1}z\,c\,z^{2^{k}-1}=z^{2^{k}},
\]
and $c^{-1}zc^{2}=cz$.  Adding the three contributions gives exactly
$\ell_{0,c}(z)$.  (We used $A^{2}=D^{2^{n}}=D$, which holds by Lemma
\ref{lem:frob}.)
\end{proof}

\begin{proposition}\label{prop:case1}
Suppose there is \emph{no} $\gamma\in L$ with $c=\gamma^{2^{k}+1}+\gamma$.  Then
$\ell_c$ has at most one root in $L$.
\end{proposition}

\begin{proof}
First, $c\neq0$: otherwise $\gamma=0$ would do.  Let $x,y$ be roots of $\ell_c$
and set $z=x+y$, so $\ell_{0,c}(z)=0$ by linearity.  If $z\neq0$, Lemma
\ref{lem:factorization} forces $A^{2^{k}+1}+A+c=0$, i.e.
$c=A^{2^{k}+1}+A$ (characteristic two), contradicting the hypothesis.  Hence
$z=0$, i.e. $x=y$.
\end{proof}

\subsection{Second case: $c=\gamma^{2^{k}+1}+\gamma$}

Assume now $c=\gamma^{2^{k}+1}+\gamma$ and $c\neq0$; then $\gamma\neq0$ and, since
$c=\gamma(\gamma^{2^{k}}+1)$, also $\gamma^{2^{k}}\neq1$, i.e. $\gamma\neq1$.
Introduce
\[
  Q(X)\;=\;c\,X^{2^{k}}+\gamma^{2}X+\gamma .
\]

\subsubsection{The auxiliary element $\Delta$}

\begin{lemma}\label{lem:delta}
There is a $\Delta\in L^{\times}$ with
\begin{equation}\label{eq:deltadef}
  \bigl(\Delta^{2^{k}-1}\bigr)^{-1}\;=\;\gamma^{2^{k}-1}+\gamma^{-1},
\end{equation}
and any such $\Delta$ satisfies
\begin{equation}\label{eq:deltafrob}
  \Delta^{2^{k}}\;=\;\gamma\Delta+(\gamma\Delta)^{2^{k}} .
\end{equation}
\end{lemma}

\begin{proof}
We have $\gamma^{2^{k}-1}+\gamma^{-1}=\gamma^{-1}(\gamma^{2^{k}}+1)\neq0$ because
$\gamma\ne 1$.  By Lemma \ref{lem:gcdpow}(1) the map $u\mapsto u^{2^{k}-1}$ is a
bijection of $L^{\times}$, so \eqref{eq:deltadef} has a solution
$\Delta\in L^{\times}$.  Rewriting \eqref{eq:deltadef} as
$\Delta^{2^{k}-1}\gamma^{-1}(\gamma^{2^{k}}+1)=1$ and multiplying by
$\gamma\Delta$ gives $\Delta^{2^{k}}(\gamma^{2^{k}}+1)=\gamma\Delta$, which is
\eqref{eq:deltafrob}.
\end{proof}

\begin{lemma}\label{lem:tracedelta}
$\Tr(\Delta)=0$.
\end{lemma}

\begin{proof}
Apply $\Tr$ to \eqref{eq:deltafrob}.  Since $\Tr$ is additive and
Frobenius-invariant, $\Tr(\Delta^{2^{k}})=\Tr(\Delta)$ and
$\Tr\bigl((\gamma\Delta)^{2^{k}}\bigr)=\Tr(\gamma\Delta)$, so
$\Tr(\Delta)=\Tr(\gamma\Delta)+\Tr(\gamma\Delta)=0$.
\end{proof}

\subsubsection{Factorisation of the zero set of $\ell_c$ through $Q$}

\begin{proposition}\label{prop:ellQ}
For all $x\in L$,
\[
  \ell_c(x)\;=\;Q(x)^{2^{k}}+\bigl(\Delta^{2^{k}-1}\bigr)^{-1}Q(x),
\]
and therefore
\begin{equation}\label{eq:F}
  \ell_c(x)=0\iff Q(x)=0\ \text{ or }\ Q(x)=\Delta^{-1}. \tag{F}
\end{equation}
\end{proposition}

\begin{proof}
Write $\lambda=\bigl(\Delta^{2^{k}-1}\bigr)^{-1}=\gamma^{2^{k}-1}+\gamma^{-1}$.
Expanding,
\[
  Q(x)^{2^{k}}=c^{2^{k}}x^{2^{2k}}+\gamma^{2^{k+1}}x^{2^{k}}+\gamma^{2^{k}},
  \qquad
  \lambda Q(x)=\lambda c\,x^{2^{k}}+\lambda\gamma^{2}x+\lambda\gamma .
\]
Now $\lambda\gamma^{2}=\gamma^{2^{k}+1}+\gamma=c$ and
$\lambda\gamma=\gamma^{2^{k}}+1$, whence the constant terms add up to
$\gamma^{2^{k}}+\gamma^{2^{k}}+1=1$ and the coefficient of $x$ is $c$.  Finally
\[
  \lambda c=(\gamma^{2^{k}-1}+\gamma^{-1})(\gamma^{2^{k}+1}+\gamma)
  =\gamma^{2^{k+1}}+\gamma^{2^{k}}+\gamma^{2^{k}}+1=\gamma^{2^{k+1}}+1 ,
\]
so the coefficient of $x^{2^{k}}$ is $\gamma^{2^{k+1}}+\gamma^{2^{k+1}}+1=1$.
Altogether $Q(x)^{2^{k}}+\lambda Q(x)=c^{2^{k}}x^{2^{2k}}+x^{2^{k}}+cx+1=\ell_c(x)$.

For \eqref{eq:F}: $\ell_c(x)=0$ means $Q(x)\bigl(Q(x)^{2^{k}-1}+\lambda\bigr)=0$,
i.e. $Q(x)=0$ or $Q(x)^{2^{k}-1}=\lambda=(\Delta^{-1})^{2^{k}-1}$.  By Lemma
\ref{lem:gcdpow}(1) the latter is equivalent to $Q(x)=\Delta^{-1}$.
\end{proof}

\subsubsection{Roots of $Q$ do not solve the Kasami equation}

\begin{lemma}\label{lem:Qroot}
Let $x$ be a root of $Q$.  Then
\begin{enumerate}
\item $x^{2^{k}}=\gamma x+(\gamma x)^{2^{k}}+1$, and more generally for all
      $i\ge0$
      \[
        x^{2^{(i+1)k}}=(\gamma x)^{2^{ik}}+(\gamma x)^{2^{(i+1)k}}+1 ;
      \]
\item $\Tr(x)=n\cdot 1$ (i.e. $\Tr(x)=n\bmod 2$ in $\F_2\subseteq L$);
\item $S(x)=\gamma x+(\gamma x)^{2}+k'\cdot1$;
\item $c\,x^{2^{k}+1}=\gamma x+(\gamma x)^{2}$.
\end{enumerate}
Consequently
\begin{equation}\label{eq:efinal}
  c\,x^{2^{k}+1}+S(x)+\alpha\Tr(x)\;=\;(k'+\alpha n)\cdot 1 .
\end{equation}
\end{lemma}

\begin{proof}
(1) $Q(x)=0$ reads $(\gamma^{2^{k}+1}+\gamma)x^{2^{k}}+\gamma^{2}x+\gamma=0$;
dividing by $\gamma\neq0$ gives
$(\gamma^{2^{k}}+1)x^{2^{k}}+\gamma x+1=0$, i.e.
$x^{2^{k}}=(\gamma x)^{2^{k}}+\gamma x+1$.  The general form follows by applying
the automorphism $u\mapsto u^{2^{ik}}$ (note $1^{2^{ik}}=1$).

(2) Apply $u\mapsto u^{2^{j}}$ to the case $i=0$ of (1) and sum over
$j=0,\dots,n-1$.  The left-hand side gives $\sum_{j}x^{2^{k+j}}=\Tr(x)$ by Lemma
\ref{lem:frob}; the right-hand side gives
$\Tr(\gamma x)+\Tr(\gamma x)+n\cdot1=n\cdot1$.

(3) Summing the identity of (1) for $i=0,\dots,k'-1$, the middle terms telescope:
\[
  S(x)=\sum_{i=0}^{k'-1}x^{2^{(i+1)k}}
      =\gamma x+(\gamma x)^{2^{k'k}}+k'\cdot 1 .
\]
By \eqref{eq:bezout}, $k'k\equiv1\pmod n$, so $(\gamma x)^{2^{k'k}}=(\gamma x)^{2}$.

(4) Multiply $x^{2^{k}}=(\gamma x)^{2^{k}}+\gamma x+1$ by $x$:
$x^{2^{k}+1}(1+\gamma^{2^{k}})=\gamma x^{2}+x$.  Multiplying by $\gamma$ and using
$c=\gamma(\gamma^{2^{k}}+1)$ gives
$c\,x^{2^{k}+1}=\gamma^{2}x^{2}+\gamma x=(\gamma x)^{2}+\gamma x$.

Finally, adding (3), (4) and $\alpha$ times (2), the terms $\gamma x$ and
$(\gamma x)^{2}$ cancel in characteristic two and \eqref{eq:efinal} results.
\end{proof}

\begin{proposition}\label{prop:noQzero}
Assume $k'+\alpha n\equiv1\pmod2$.  Then no root of $Q$ satisfies \eqref{eq:K}.
\end{proposition}

\begin{proof}
If $x$ satisfies \eqref{eq:K} then $c\,x^{2^{k}+1}+S(x)+\alpha\Tr(x)=0$, whereas
if $Q(x)=0$ then by \eqref{eq:efinal} this same quantity equals
$(k'+\alpha n)\cdot1=1$.  Since $1\neq0$ in $L$, the two are incompatible.
\end{proof}

\subsubsection{Uniqueness inside the fibre $Q(x)=\Delta^{-1}$}

\begin{lemma}\label{lem:pairsum}
Let $x_0\neq x_1$ satisfy $Q(x_0)=Q(x_1)=\Delta^{-1}$.  Then
\[
  x_0+x_1=\Delta,\qquad \Tr(x_0)=\Tr(x_1),\qquad S(x_0)+S(x_1)=\gamma\Delta+(\gamma\Delta)^{2}.
\]
\end{lemma}

\begin{proof}
We first rewrite the fibre condition.  Dividing $Q(x)=\Delta^{-1}$, i.e.
$c\,x^{2^{k}}+\gamma^{2}x+\gamma=\Delta^{-1}$, by $\gamma^{2}\Delta$ and using
\[
  \frac{c}{\gamma^{2}\Delta}=\frac{\gamma^{2^{k}}+1}{\gamma\Delta}
  =\frac{1}{\Delta^{2^{k}}}
\]
(the last equality being \eqref{eq:deltafrob} in the form
$\Delta^{2^{k}}(\gamma^{2^{k}}+1)=\gamma\Delta$), we find that
$Q(x)=\Delta^{-1}$ is equivalent to
\[
  u+u^{2^{k}}=(\gamma\Delta)^{-1}+\bigl((\gamma\Delta)^{2}\bigr)^{-1},
  \qquad u=x/\Delta .
\]
Thus $u_0=x_0/\Delta$ and $u_1=x_1/\Delta$ have the same image under the additive
map $u\mapsto u^{2^{k}}+u$, whose kernel is $\{0,1\}$ by Lemma
\ref{lem:gcdpow}(3).  So $u_0=u_1$ or $u_0=u_1+1$; as $x_0\neq x_1$ the second
holds, i.e. $x_0+x_1=\Delta$.

By Lemma \ref{lem:tracedelta} and additivity, $\Tr(x_0)+\Tr(x_1)=\Tr(\Delta)=0$.

For the last identity, additivity of $S$ gives $S(x_0)+S(x_1)=S(\Delta)$.  Applying
$u\mapsto u^{2^{ik}}$ to \eqref{eq:deltafrob} yields
$\Delta^{2^{(i+1)k}}=(\gamma\Delta)^{2^{ik}}+(\gamma\Delta)^{2^{(i+1)k}}$, and
summing over $i=0,\dots,k'-1$ telescopes to
\[
  S(\Delta)=\gamma\Delta+(\gamma\Delta)^{2^{k'k}}=\gamma\Delta+(\gamma\Delta)^{2},
\]
using $k'k\equiv1\pmod n$ once more.  (In contrast with Lemma \ref{lem:Qroot}(3)
there is no constant term here, since \eqref{eq:deltafrob} is homogeneous.)
\end{proof}

For $x\in L$ set
\[
  e(x)\;=\;c\,x^{2^{k}+1}+S(x)+\alpha\Tr(x),
\]
so that, by Lemma \ref{lem:fibre}, $x$ solves \eqref{eq:K} exactly when $e(x)=0$.

\begin{proposition}\label{prop:e0e1}
Let $x_0\neq x_1$ satisfy $Q(x_0)=Q(x_1)=\Delta^{-1}$.  Then
$e(x_0)+e(x_1)=1$.
\end{proposition}

\begin{proof}
By Lemma \ref{lem:pairsum} we may write $x_1=x_0+\Delta$, and
\[
  e(x_0)+e(x_1)=c\bigl(x_0^{2^{k}+1}+x_1^{2^{k}+1}\bigr)+\gamma\Delta+(\gamma\Delta)^{2}.
\]
Expanding $x_1^{2^{k}+1}=(x_0+\Delta)^{2^{k}}(x_0+\Delta)$ in characteristic two,
\[
  x_0^{2^{k}+1}+x_1^{2^{k}+1}
   =\Delta\,x_0^{2^{k}}+\Delta^{2^{k}}x_0+\Delta^{2^{k}+1}
   =\Delta\,x_0^{2^{k}}+\Delta^{2^{k}}\bigl(x_0+\Delta\bigr).
\]
From $Q(x_0)=\Delta^{-1}$ we get $c\,x_0^{2^{k}}=\Delta^{-1}+\gamma^{2}x_0+\gamma$,
hence
\[
  c\,\Delta\,x_0^{2^{k}}=1+\gamma^{2}\Delta x_0+\gamma\Delta .
\]
From \eqref{eq:deltafrob}, $\Delta^{2^{k}}(\gamma^{2^{k}}+1)=\gamma\Delta$, and
since $c=\gamma(\gamma^{2^{k}}+1)$ this gives
$c\,\Delta^{2^{k}}=\gamma\cdot\gamma\Delta=\gamma^{2}\Delta$.  Therefore
\[
  c\bigl(x_0^{2^{k}+1}+x_1^{2^{k}+1}\bigr)
  =1+\gamma^{2}\Delta x_0+\gamma\Delta+\gamma^{2}\Delta\bigl(x_0+\Delta\bigr)
  =1+\gamma\Delta+(\gamma\Delta)^{2},
\]
and adding $\gamma\Delta+(\gamma\Delta)^{2}$ leaves $e(x_0)+e(x_1)=1$.
\end{proof}

\begin{proposition}\label{prop:case2}
Assume $k'+\alpha n\equiv1\pmod2$ and $c=\gamma^{2^{k}+1}+\gamma$ for some
$\gamma\in L$.  Then \eqref{eq:K} has at most one nonzero solution.
\end{proposition}

\begin{proof}
If $c=0$ then \eqref{eq:K} reads $\rhs(x)=0$, and Proposition
\ref{prop:zerofibre} gives $x=0$; so there is no nonzero solution at all.
Assume $c\ne 0$; then $\gamma\notin\{0,1\}$ as observed above, and $\Delta$ is
available.  Let $x_0\neq x_1$ be nonzero solutions of \eqref{eq:K}.  By
Proposition \ref{prop:linearized} both are roots of $\ell_c$, so by \eqref{eq:F}
each satisfies $Q(x_j)=0$ or $Q(x_j)=\Delta^{-1}$.  Proposition
\ref{prop:noQzero} excludes the first alternative, so
$Q(x_0)=Q(x_1)=\Delta^{-1}$.  Proposition \ref{prop:e0e1} then gives
$e(x_0)+e(x_1)=1$, contradicting $e(x_0)=e(x_1)=0$.
\end{proof}

\subsection{Proof of Theorem \ref{thm:kasami}}

\begin{proposition}\label{prop:unique}
Assume $k'+\alpha n\equiv1\pmod2$.  Let $c,x,y\in L$ satisfy \eqref{eq:K} for $x$
and for $y$, and suppose that $x\neq0$ and $y\neq0$ whenever $c\neq0$.  Then
$x=y$.
\end{proposition}

\begin{proof}
If $c=0$, then $\rhs(x)=\rhs(y)=0$ and Proposition \ref{prop:zerofibre} gives
$x=y=0$.  Let $c\neq0$, so $x,y\neq0$.  If $c$ is not of the form
$\gamma^{2^{k}+1}+\gamma$, then by Proposition \ref{prop:linearized} both $x$ and
$y$ are roots of $\ell_c$, and Proposition \ref{prop:case1} gives $x=y$.  If $c$
is of that form, Proposition \ref{prop:case2} gives $x=y$.
\end{proof}

\begin{proof}[Proof of Theorem \ref{thm:kasami}]
Necessity is Proposition \ref{prop:necessity}.  For sufficiency assume
$k'+\alpha n\equiv1\pmod2$ and let $q_\alpha(x)=q_\alpha(y)=:c$.

If $c=0$ then by Lemma \ref{lem:fibre}(2) $\rhs(x)=\rhs(y)=0$, so $x=y=0$ by
Proposition \ref{prop:zerofibre}.  If $c\neq0$ then $x\neq0\neq y$ (as
$q_\alpha(0)=0$), and by Lemma \ref{lem:fibre}(1) both satisfy \eqref{eq:K};
Proposition \ref{prop:unique} yields $x=y$.  Hence the evaluation map of
$q_\alpha$ is injective, and since $L$ is finite it is bijective.  Thus
$q_\alpha$ is a permutation polynomial.
\end{proof}

\section{The transfer theorem}\label{sec:transfer}

We now relate $q_\alpha$ and $P_\beta$ through two linearized maps.  Fix
$k',n'$ as in \eqref{eq:bezout} and $\alpha,\beta\in\{0,1\}$, and define
$L\to L$ by
\[
  \psi_\beta(x)=\sum_{i=0}^{k-1}x^{2^{i}}+\beta\Tr(x),
  \qquad
  \varphi_\alpha(x)=\sum_{i=0}^{k'-1}x^{2^{ik}}+\alpha\Tr(x) .
\]
Both are $\F_2$-linear.

\begin{lemma}\label{lem:Rphi}
For all $x\in L$, $\rhs(x)=\varphi_\alpha(x)^{2^{k}}$.
\end{lemma}

\begin{proof}
$\bigl(\sum_{i=0}^{k'-1}x^{2^{ik}}\bigr)^{2^{k}}=\sum_{i=1}^{k'}x^{2^{ik}}=S(x)$
and $\Tr(x)^{2^{k}}=\Tr(x)$, while $\alpha^{2^{k}}=\alpha$.
\end{proof}

\begin{lemma}\label{lem:Pdiv}
For $x\neq0$, $P_\beta(x)=\psi_\beta(x)^{2^{k}+1}\big/x^{2^{k}}$.
\end{lemma}

\begin{proof}
$P_\beta(x)=\psi_\beta(x)^{2^{k}+1}x^{(2^{n}-1)-2^{k}}$ by definition (note
$2^{k}\le 2^{n}-1$ since $k<n$), and $x^{(2^{n}-1)-2^{k}}=x^{2^{n}-1}x^{-2^{k}}
=x^{-2^{k}}$.
\end{proof}

\begin{proposition}\label{prop:transferinv}
Assume \eqref{eq:bezout}, $k'+\alpha n\equiv 1\pmod2$ and
$\beta\equiv n'+\alpha k\pmod 2$.  Then $\varphi_\alpha\circ\psi_\beta=\mathrm{id}_L$;
consequently $\psi_\beta$ and $\varphi_\alpha$ are mutually inverse bijections of $L$.
\end{proposition}

\begin{proof}
Fix $x$ and abbreviate $T=\Tr(x)$.  First,
\[
  \sum_{i=0}^{k'-1}\Bigl(\sum_{j=0}^{k-1}x^{2^{j}}\Bigr)^{2^{ik}}
  =\sum_{i=0}^{k'-1}\sum_{j=0}^{k-1}x^{2^{ik+j}}
  =\sum_{m=0}^{k'k-1}x^{2^{m}}
  =\sum_{m=0}^{n'n}x^{2^{m}}
  =x+n'\,T ,
\]
where the last step splits the range $\{0,\dots,n'n\}$ into the singleton
$\{n'n\}$ contributing $x^{2^{n'n}}=x$ (Lemma \ref{lem:frob}) and $n'$ blocks of
$n$ consecutive exponents, each contributing $T$.  Second,
$\bigl(\beta T\bigr)^{2^{ik}}=\beta T$, so the $\beta$-part of
$\varphi_\alpha(\psi_\beta(x))$ contributes $k'\beta T$.  Third,
\[
  \alpha\Tr\bigl(\psi_\beta(x)\bigr)
  =\alpha\Bigl(\sum_{j=0}^{n-1}\sum_{i=0}^{k-1}x^{2^{i+j}}+\beta\sum_{j=0}^{n-1}T\Bigr)
  =\alpha\,(k+\beta n)\,T .
\]
Altogether
\[
  \varphi_\alpha\bigl(\psi_\beta(x)\bigr)=x+\bigl(n'+\beta k'+\alpha k+\alpha\beta n\bigr)T ,
\]
and the integer coefficient is to be read modulo $2$.  Using
$\beta\equiv n'+\alpha k$ and $k'+\alpha n\equiv1$,
\[
  n'+\beta k'+\alpha k+\alpha\beta n=(n'+\alpha k)+\beta(k'+\alpha n)
  \equiv\beta+\beta\equiv0 \pmod 2 .
\]
Hence $\varphi_\alpha\circ\psi_\beta=\mathrm{id}$.  Since $L$ is finite, both maps
are bijections and each is the inverse of the other.
\end{proof}

\begin{theorem}[Transfer theorem; Dobbertin, Thm.~4]\label{thm:transfer}
Assume $0<k<n$, $\gcd(k,n)=1$, $0<k'<n$, \eqref{eq:bezout},
$k'+\alpha n\equiv1\pmod2$ and $\beta\equiv n'+\alpha k\pmod2$.  Then
\begin{enumerate}
\item $\psi_\beta$ and $\varphi_\alpha$ are mutually inverse bijections of $L$;
\item $q_\alpha\circ\psi_\beta=(\,\cdot\,)^{-1}\circ P_\beta$ and
      $P_\beta\circ\varphi_\alpha=(\,\cdot\,)^{-1}\circ q_\alpha$;
\item $q_\alpha$ and $P_\beta$ are both permutation polynomials of $L$.
\end{enumerate}
\end{theorem}

\begin{proof}
(1) is Proposition \ref{prop:transferinv}.

(2) Let $x\in L$.  If $x=0$ then $\psi_\beta(0)=0$, and $q_\alpha(0)=0=P_\beta(0)^{-1}$
(recall $0^{-1}=0$), so the first identity holds.  Let $x\neq0$ and put
$y=\psi_\beta(x)$; then $y\neq0$, since $\varphi_\alpha(y)=x\neq0=\varphi_\alpha(0)$.
By Lemma \ref{lem:fibre}(1) applied to $c=q_\alpha(y)$, and by Lemma
\ref{lem:Rphi},
\[
  q_\alpha(y)\,y^{2^{k}+1}=\rhs(y)=\varphi_\alpha(y)^{2^{k}}=x^{2^{k}} .
\]
By Lemma \ref{lem:Pdiv}, $P_\beta(x)=y^{2^{k}+1}/x^{2^{k}}\neq0$, so
$q_\alpha(y)\,P_\beta(x)=1$, i.e. $q_\alpha(\psi_\beta(x))=P_\beta(x)^{-1}$.  The
second identity follows by substituting $x\mapsto\varphi_\alpha(x)$ and using
$\psi_\beta\circ\varphi_\alpha=\mathrm{id}$.

(3) $q_\alpha$ is a permutation polynomial by Theorem \ref{thm:kasami}, the parity
hypothesis being satisfied.  Hence $(\,\cdot\,)^{-1}\circ q_\alpha$ is a bijection
of $L$, so by (2) $P_\beta\circ\varphi_\alpha$ is a bijection; since
$\varphi_\alpha$ is a bijection, so is $P_\beta$.
\end{proof}

\section{The MCM polynomial}\label{sec:mcm}

\begin{lemma}\label{lem:mcmid}
For $0<k<n$, the polynomials $M_{n,k}$ and $P_0$ induce the same evaluation map on
$L$.
\end{lemma}

\begin{proof}
At $x=0$ both vanish: $M_{n,k}(0)=0$ because of the factor $z$, and
$P_0(0)=0$ because $\sum_{i=0}^{k-1}0^{2^{i}}=0$ (all exponents $2^{i}$ are
positive) and $k\ge1$.

Let $x\neq0$, so $x^{2^{n}-1}=1$.  Since $2^{i}-1\ge0$,
\[
  \sum_{i=0}^{k-1}x^{2^{i}}=x\sum_{i=0}^{k-1}x^{2^{i}-1}=:x\,\Sigma .
\]
Therefore, using $2^{k}\le2^{n}-1$,
\[
  P_0(x)=\bigl(x\Sigma\bigr)^{2^{k}+1}x^{(2^{n}-1)-2^{k}}
        =\Sigma^{2^{k}+1}\,x^{2^{k}+1+(2^{n}-1)-2^{k}}
        =\Sigma^{2^{k}+1}\,x\,x^{2^{n}-1}
        =x\,\Sigma^{2^{k}+1}=M_{n,k}(x). \qedhere
\]
\end{proof}

\begin{proof}[Proof of Theorem \ref{thm:main}]
Let $0<k<n$, $\gcd(k,n)=1$ and $k$ odd.  Choose $k'$ with $0<k'<n$ and
$kk'\equiv1\pmod n$, and let $n'$ be as in \eqref{eq:bezout}, i.e.
$k'k=n'n+1$.  Apply Theorem \ref{thm:transfer} with
\[
  \alpha\equiv n' \pmod 2,\qquad \beta=0 .
\]
Its two parity hypotheses hold:
\begin{itemize}
\item Reducing \eqref{eq:bezout} modulo $2$ and using that $k$ is odd gives
      $k'\equiv k'k\equiv n'n+1\pmod2$; hence
      \[
        k'+\alpha n\equiv (n'n+1)+n'n\equiv1 \pmod 2 .
      \]
\item $\beta=0$ and, again since $k$ is odd,
      $n'+\alpha k\equiv n'+n'k=n'(1+k)\equiv0\pmod 2$, so
      $\beta\equiv n'+\alpha k \pmod 2$.
\end{itemize}
Theorem \ref{thm:transfer}(3) therefore shows that $P_0$ is a permutation
polynomial of $L$, and by Lemma \ref{lem:mcmid} the evaluation map of $M_{n,k}$
coincides with that of $P_0$.  Hence $M_{n,k}$ is a permutation polynomial of
$\GF(2^{n})$.
\end{proof}

\begin{remark}
The hypothesis that $k$ is odd is used exactly twice, namely in the two parity
verifications above; it is what makes the choice $\beta=0$ (the ``classical'' MCM
polynomial, without trace term) compatible with the permutation parity
$k'+\alpha n\equiv1$ of the Kasami side.
\end{remark}

\appendix

\section{Correspondence with the formal development}\label{app:lean}

Every statement of this note corresponds to a machine-checked declaration of the
accompanying Lean development (namespace \texttt{GeneralizedKasami}).  In the Lean
files the field is \texttt{GF n}, the parameters $\alpha,\beta$ are elements of
\texttt{Fin 2} named \texttt{a}, \texttt{b}, and $k'$, $n'$, $\Delta$ are
\texttt{kInv}, \texttt{nPrime}, \texttt{Delta}.

\begin{center}
\begin{tabular}{lll}
\hline
Here & Lean declaration & File\\
\hline
$\Tr$ & \texttt{tracePolynomial} & \texttt{basic.lean}\\
$q_\alpha$, its evaluation & \texttt{genKasamiPol}, \texttt{genKasamiEval} & \texttt{basic.lean}\\
$\rhs$ & \texttt{equationRHS} & \texttt{basic.lean}\\
$\ell_c$, $\ell_{0,c}$ & \texttt{ell}, \texttt{ell\_0} & \texttt{basic.lean}\\
$Q$, $c=\gamma^{2^{k}+1}+\gamma$ & \texttt{gammaPolynomial}, \texttt{gammaCoefficient} & \texttt{GammaSteps.lean}\\
$S$ & \texttt{iteratedKasamiSum} & \texttt{GammaSteps.lean}\\
$\Delta$ & \texttt{IsDelta} & \texttt{DeltaSteps.lean}\\
$e(x)$ & \texttt{eTerm} & \texttt{PairSteps.lean}\\
$P_\beta$, $\psi_\beta$, $\varphi_\alpha$ & \texttt{genMCMpol}, \texttt{psi\_beta}, \texttt{phi\_alpha} & \texttt{genMCM.lean}\\
$M_{n,k}$ & \texttt{MCMpolynomial}, \texttt{MCMeval} & \texttt{MCMCorollary.lean}\\
\hline
Lemma \ref{lem:frob} & \texttt{gf\_frobenius\_pow\_card}, \texttt{gf\_frobenius\_mod} & \texttt{basic.lean}\\
Lemma \ref{lem:fibre} & \texttt{evaluation\_eq\_iff\_equation\_of\_ne\_zero}, & \texttt{basic.lean}\\
 & \texttt{evaluation\_eq\_zero\_iff\_equation} & \\
Prop. \ref{prop:necessity} & \texttt{permutation\_implies\_parity} & \texttt{basic.lean}\\
Lemma \ref{lem:shift} & \texttt{kasami\_sum\_frobenius}, \texttt{trace\_sum\_frobenius} & \texttt{basic.lean}\\
Prop. \ref{prop:linearized} & \texttt{kasamiEquation\_implies\_linearized} & \texttt{basic.lean}\\
Prop. \ref{prop:zerofibre} & \texttt{equationRHS\_eq\_zero\_unique} & \texttt{basic.lean}\\
Lemma \ref{lem:factorization} & \texttt{ell\_0\_eval\_factorization} & \texttt{basic.lean}\\
Prop. \ref{prop:case1} & \texttt{linearized\_unique\_of\_not\_exists\_gamma} & \texttt{basic.lean}\\
Lemma \ref{lem:delta} & \texttt{exists\_delta}, \texttt{delta\_frobenius\_identity} & \texttt{DeltaSteps.lean}\\
Lemma \ref{lem:tracedelta} & \texttt{trace\_sum\_delta\_eq\_zero} & \texttt{DeltaSteps.lean}\\
Prop. \ref{prop:ellQ} & \texttt{ell\_eval\_eq\_gammaPolynomial\_frobenius}, & \texttt{EllSteps.lean}\\
 & \texttt{ell\_eval\_eq\_zero\_iff\_gammaPolynomial} & \\
Lemma \ref{lem:Qroot} & \texttt{gammaPolynomial\_eq\_zero\_iff(\_iterate)}, & \texttt{GammaSteps.lean}\\
 & \texttt{trace\_sum\_of\_gammaPolynomial\_root}, & \\
 & \texttt{iteratedKasamiSum\_of\_gammaPolynomial\_root}, & \\
 & \texttt{gammaPolynomial\_root\_final\_identity} & \\
Prop. \ref{prop:noQzero} & \texttt{kasamiEquation\_not\_gammaPolynomial\_zero} & \texttt{genkasamipol.lean}\\
Lemma \ref{lem:pairsum} & \texttt{gammaPolynomial\_eq\_delta\_inv\_iff}, & \texttt{DeltaSteps.lean}\\
 & \texttt{frobenius\_add\_eq\_frobenius\_add\_iff}, & \\
 & \texttt{pair\_sum\_eq\_delta}, \texttt{pair\_trace\_eq}, & \texttt{PairSteps.lean}\\
 & \texttt{iteratedKasamiSum\_delta} & \\
Prop. \ref{prop:e0e1} & \texttt{e0\_plus\_e1\_equals\_one} & \texttt{PairSteps.lean}\\
Prop. \ref{prop:case2} & \texttt{kasamiEquation\_unique\_of\_exists\_gamma} & \texttt{genkasamipol.lean}\\
Prop. \ref{prop:unique} & \texttt{kasamiEquation\_unique} & \texttt{genkasamipol.lean}\\
Theorem \ref{thm:kasami} & \texttt{parity\_implies\_permutation}, & \texttt{genkasamipol.lean}\\
 & \texttt{generalizedKasami\_isPermutation\_iff} & \\
Lemma \ref{lem:Rphi} & \texttt{equationRHS\_eq\_phi\_pow} & \texttt{genMCM.lean}\\
Lemma \ref{lem:Pdiv} & \texttt{genMCMeval\_eq\_div} & \texttt{genMCM.lean}\\
Prop. \ref{prop:transferinv} & \texttt{transfer\_maps\_inverse} & \texttt{genMCM.lean}\\
Theorem \ref{thm:transfer} & \texttt{dobbertin\_theorem4} & \texttt{genMCM.lean}\\
Lemma \ref{lem:mcmid} & \texttt{MCMeval\_eq\_genMCMeval\_zero} & \texttt{MCMCorollary.lean}\\
Theorem \ref{thm:main} & \texttt{MCMpolynomial\_isPermutation} & \texttt{MCMCorollary.lean}\\
\hline
\end{tabular}
\end{center}

\medskip
\noindent
The Bézout data \eqref{eq:bezout} is produced in Lean by
\texttt{exists\_positive\_inverse\_coefficients}
(\texttt{MCMCorollary.lean}).  The formal statement of Theorem \ref{thm:main} is
\begin{verbatim}
theorem MCMpolynomial_isPermutation
    (n k : Nat) (hn : 0 < n) (hk : 0 < k) (hkn : k < n)
    (hcop : Nat.Coprime k n) (hkodd : Odd k) :
    IsPermutationPolynomial (MCMpolynomial n k)
\end{verbatim}
(the natural-number type is written \texttt{Nat} here for typesetting reasons),
where \texttt{IsPermutationPolynomial p} unfolds to \texttt{Function.Bijective p.eval}.

\end{document}
